% ================================
% === General Response ===
% ================================
\section{General Response}

We thank the Associate Editor and all reviewers for the careful second-round assessment. Three concerns recurred across the reports. These were the theoretical basis of the functional decoupling, the absence of statistical validation, and evaluation on a single benchmark. Each is answered with new analysis or new experiments in the manuscript. The table below maps every concern to the change and its location.

\begin{changetable}{p{0.24\linewidth}p{0.42\linewidth}p{0.25\linewidth}}
\textbf{Concern} & \textbf{How we address it} & \textbf{Where} \\
\hline
Theory of the decoupling (R2, R4) & Exact entropy identity, one derivation per module, and a $2\times2$ ambiguity-stratified analysis testing their prediction & Sec.~\ref{sec:decoupling_principle}, and Sec.~\ref{sec:quadrants} with Table~\ref{tab:quadrants} \\
Statistical validation (R2, R6) & Five seeds, $95\%$ $t$-intervals, paired $t$-tests against every matched-protocol comparator with Holm correction & Sec.~\ref{sec:sota_comparison} and Table~\ref{tab:main_results}, with full statistics in Table~\ref{tab:significance_letter} of this letter \\
Single benchmark, dataset scale (R2, R4) & Cross-dataset study on a larger EMOTIC benchmark, plus the Flickr-LDL and Emotion6 label-distribution benchmarks & Sec.~\ref{sec:crossdataset}, Tables~\ref{tab:emotic_pool} and~\ref{tab:ldl_transfer} \\
Recent baselines (R4, EiC) & Two 2026 intent methods added and compared, one under a matched backbone, and four recent MLLMs evaluated zero-shot & Sec.~\ref{sec:related_work}, Tables~\ref{tab:main_results}, \ref{tab:zeroshot_vlm}, \ref{tab:backbone_ablation} \\
Abstract comparability (AE) & More intuitive presentation of advantages & Abstract \\
Proofreading, bibliography (R4, EiC) & Venue fields corrected, new works discussed individually, \res{52} items & Bibliography, Sec.~\ref{sec:related_work} \\
\end{changetable}

Two points of context apply throughout, so we state them once here.

\textbf{The cross-dataset study.} To the best of our knowledge, no larger-scale visual intention recognition dataset currently matches our target in terms of ambiguous recognition and subjective annotations. We therefore selected three closely related affective understanding benchmarks to evaluate generalization under complementary settings, covering multi-label emotion recognition and label-distribution learning.
It spans two axes, and both are in the manuscript. EMOTIC is the larger benchmark ($23{,}571$ images, $34{,}320$ annotated instances), and it keeps our multi-label output geometry. It releases multi-annotator labels only for validation and test, so we report a controlled cross-validation on that multi-annotated pool rather than a leaderboard comparison. Flickr-LDL and Emotion6 are smaller, but carry native subjective label distributions under standard published protocols, and remove threshold selection entirely. Together the two axes answer the scale question and the breadth question.

\textbf{What this letter adds.} The manuscript is at the $35$-page limit, so three items live here instead. These are the complete per-method significance statistics (Table~\ref{tab:significance_letter}), the seed dispersion behind the cross-dataset table (Table~\ref{tab:ldl_letter}), and the per-class F1 attribution (Table~\ref{tab:per_class_attr_letter}) behind the manuscript's one-sentence interpretability summary. Everything else described below is in the revised manuscript, which is marked in red.
